\section{Project Plan}

To ensure the systematic development and evaluation of the proposed dual-objective 
framework within the allotted research timeframe, the project will be executed using 
structured work packages, defined milestones, and proactive risk management protocols.

\subsection{Key Tasks}
The research methodology is partitioned into five sequential and logically dependent 
tasks.

\begin{enumerate}

\item \textbf{Task 1: Infrastructure and Environment Provisioning}

\textbf{Objectives:} Establish a deterministic, hardware-agnostic execution environment.

\textbf{Tasks:}
\begin{itemize}
\item Containerize the development stack using Docker and Nix.
\item Deploy the base open-weights model (e.g., Llama-3-8B) on the target GPU 
infrastructure utilizing the vLLM serving engine.
\item Establish baseline latency metrics for unconstrained generation.
\end{itemize}

\item \textbf{Task 2: Schema Compiler and Syntax Oracle Engineering}

\textbf{Objectives:} Develop the C++ backend for the symbolic constraint engine.

\textbf{Tasks:}
\begin{itemize}
\item Construct the Schema Compiler to parse annotated JSON Schemas.
\item Implement the Finite State Machine (FSM) generator for the Context-Free 
Grammar (CFG) mask.
\item Develop the topological-sort engine for evaluating the Cross-Field 
Dependency Graph ($G$).
\item Compile the \texttt{pybind11} bridges.
\end{itemize}

\item \textbf{Task 3: Cascaded Verification and KV-Cache Rollback}

\textbf{Objectives:} Integrate the Syntax Oracle into the autoregressive decoding loop.

\textbf{Tasks:}
\begin{itemize}
\item Subclass the \texttt{LogitsProcessor} to inject the token-level $O(1)$ FSM mask.
\item Implement the element-level trigger hooks.
\item Engineer the speculative KV-cache rollback mechanism to 
truncate and regenerate upon Context-Sensitive Invariant (CSI) violations.
\end{itemize}

\item \textbf{Task 4: Data Synthesis and Semantic Internalization}

\textbf{Objectives:} Align the neural policy to the target Domain-Specific Language (DSL).

\textbf{Tasks:}
\begin{itemize}
\item Utilize the Oracle to run coverage-guided random walks, 
synthesizing a mathematically perfect dataset $D$.
\item Execute Supervised Fine-Tuning (SFT) using Low-Rank 
Adaptation (LoRA) with optimized rank capacity ($r=16, \alpha=32$).
\end{itemize}

\item \textbf{Task 5: Evaluation, Benchmarking, and Writing}

\textbf{Objectives:} Quantify the architectural synergy and draft the final dissertation.

\textbf{Tasks:}
\begin{itemize}
\item Measure the Entropy-Normalized Oracle Intervention Rate (OIR) 
across unaligned versus LoRA-aligned models.
\item Benchmark system throughput (tokens/second).
\item Finalize the manuscript/thesis.
\end{itemize}
\end{enumerate}

\subsection{Project Constraints}
To remain feasible, the framework will be evaluated against a singular, 
highly restricted subset of a production DSL (e.g., OpenAPI configuration). The 
compiler will support a targeted 
list of semantic annotations (e.g., arithmetic bounds, referential equality) 
rather than Turing-complete logic.

The training phase is computationally bounded by available GPU VRAM. 
The use of PEFT (LoRA) specifically mitigates this constraint, allowing 
adaptation to occur on a single consumer-grade or via standard university SLURM tasks.